\chapter{DOGMA regulation and selective transformations}
Chapter 18 exposed a collision: AACA and ACAA have identical dinucleotide
histograms, length and final base. Any output head receiving only that state
must give them the same answer. Can an ordered, selectively updated memory
separate them without attention? We will construct one such operator, prove
what it preserves, differentiate its gates, and then show why these successes
still do not establish a useful genomic model.

DOGMA remains the book's proposed non-Transformer DNA-native research direction;
Hermon DNA is the Transformer-based direction, and Evolutor is the broader
framework. The scalar and vector operators here are explicit research references,
not a released foundation model. Their merit is that their assumptions can be
tested. Their biological relevance remains an empirical question.

\section{A regulator chooses coefficients, not an explanation}
Order the 16 pairs in $\{A,C,G,T\}^2$ lexicographically in that base order.
For pair index $i$, store candidate vector $v_i\in\mathbb R^D$, write logit
$u_i\in\mathbb R^D$, and retention logit $r_i\in\mathbb R^D$. These are three
parameter tables of shape $16\times D$. The recurrent state is
$h\in\mathbb R^D$. All coordinates are dimensionless software quantities,
not concentrations or binding energies.

At position $t$, the current base and the previous valid base select $i_t$.
There is no pair at the first base. Selection is a deterministic address lookup,
not a learned top-$k$ expert decision. Chapter 10 owns that different routing
problem. This chapter asks what transformation the selected row performs.

\Plate{01-address}{A pair selects corresponding rows in three separate
parameter tables. The row index is shared, but candidate values, write gates
and retention gates have different meanings. This is software addressing;
the diagram does not represent complementary molecules binding together.}{fig:address}
\Listing{pair_ids}
The returned previous base is part of the streaming contract. To continue a
record, carry it into the next compilation call alongside the recurrent state.
To start an independent record, clear both. Carrying one without the other
creates a different model. The example accepts only canonical bases; an unknown
biological base is not silently deleted as padding.

\section{Derive a bounded selective transformation}
The logistic function is $\sigma(z)=1/(1+e^{-z})$. Define coordinatewise gates
\begin{equation}
 g_t=\sigma(u_{i_t}),\qquad \rho_t=\sigma(r_{i_t}),
 \qquad a_t=1-g_t\odot(1-\rho_t).
 \label{eq:coefficients}
\end{equation}
The symbol $\odot$ denotes coordinatewise multiplication. Let
$\widetilde h_t=\rho_t\odot h_{t-1}+(1-\rho_t)\odot v_{i_t}$ be a
candidate state transformation. Interpolate between keeping the old state
and applying that transformation:
\begin{align}
 h_t&=(1-g_t)\odot h_{t-1}+g_t\odot\widetilde h_t\\
    &=a_t\odot h_{t-1}+(1-a_t)\odot v_{i_t}.
 \label{eq:update}
\end{align}
This derivation matters: a write gate does not merely multiply the candidate.
It changes the weight of the retained state too. Using $a_t h_{t-1}+g_tv_i$
would define a different, potentially accumulating update with different bounds.
For the first base or a padding slot, define the update as exact identity;
an implementation should not approximate identity using a very negative logit.

\Plate{02-update}{The two coefficients sum to one per channel. At
$g=\rho=1/2$, $a=3/4$; old state $0.2$ and candidate $1$ produce $0.4$.
The two paths are arithmetic dependencies, not regulatory proteins.}{fig:update}
\Listing{coefficients}

\begin{proposition}[Coordinatewise range preservation]
If each initial coordinate and every candidate coordinate lie in $[-B,B]$
for finite $B$, then every state coordinate also lies in $[-B,B]$.
\end{proposition}
\begin{proof}
Since $0\leq g,\rho\leq1$, we have $0\leq a\leq1$. Each new coordinate
is a convex combination of its old value and the corresponding candidate.
The interval is convex, so induction proves the assertion. Identity slots
preserve the same interval.
\end{proof}
For fixed finite tables such a $B$ exists, but training unconstrained candidates
does not preserve a chosen global bound across optimizer steps. One could
parameterize candidates as $B\tanh(z)$; doing so introduces an additional
derivative and a modeling choice. Bounded states also do not imply that all
histories are distinguishable or that numerical gradients are well conditioned.

\section{Use the lost-order example as a constructive test}
Choose $D=1$, $v_{AA}=0$, $v_{AC}=1$, $v_{CA}=-1$, and all other candidates
zero. Set all logits to zero, so $g=\rho=1/2$ and $a=3/4$. Begin at $h_0=0$.
AACA selects AA, AC, CA. Its states after pairs are
$0,\ 1/4,\ -1/16$. ACAA selects AC, CA, AA and yields
$1/4,\ -1/16,\ -3/64$. The final difference is $-1/64$.

\begin{center}\input{\Generated/trace.tex}\end{center}
Both strings still have the same pair counts. The difference arises because
the update is ordered and earlier candidates are attenuated by later retention.
It is not evidence that training discovered a motif or that these strings
have different biological functions. A useful counterexample fixes the failure
being repaired and states exactly which extra mechanism repairs it.

\Plate{05-ablation}{Generated traces separate two inputs that collide under
the Chapter 18 histogram. Index zero is the initial state and index one is the
first base, which creates no pair. The figure demonstrates order sensitivity
of one assigned operator, not task accuracy or an optimized architecture.}{fig:ablation}

This update also differs from Chapter 18 at nonunit retention: that reference
adds a unit count, whereas Equation~\ref{eq:update} interpolates toward a
candidate. Its mass is therefore not a histogram count. Reusing the same state
dimension or DNA address table does not make the state semantics interchangeable.

\section{Implement one record with explicit ownership}
Read the implementation in three stages: validate the whole request, prepare
the input-addressed coefficients, then fold the selected transformations over
the incoming state. The validation is outside the time loop. It prevents a
bad late address from masquerading as a partially successful record; it is
not intended as a high-throughput GPU validation strategy.
\Listing{selective_scan}
The scanner returns a terminal tensor and a history that includes the initial
state. No in-place write alters a caller-supplied state. Table shapes, dtype,
device, finite values and pair IDs are validated before processing. A malformed
later ID is rejected before any state update, avoiding partially applied work.
The ID $-1$ denotes an exact no-pair operation. A padded storage slot may use
the same identity, but the token-to-pair compiler must not advance its previous
base for padding. Biological uncertainty requires a separate policy.

The scanner operates on compiled IDs rather than raw strings. This separation
makes the state update easy to test, but it also exposes a responsibility:
the caller must retain the previous base when compiling chunks. A production
state object should bundle that field, the recurrent tensor, record identity,
model revision and any position-dependent metadata. The teaching function
does not supply a serving runtime merely because it accepts a state argument.

\section{Prove causality and chunk equivalence}
The pair at $t$ depends only on bases up to $t$. The three selected table rows
are functions of that pair and fixed model parameters. By induction,
$h_t$ depends only on the valid prefix through $t$ and the initial state.
If $h_t$ predicts the next base, that is a causal next-token feature. If the
target is the current base, feeding that same base into the feature before
prediction is target leakage; causality alone does not settle label alignment.

\Plate{03-causality}{Gate coefficients can be prepared from known token pairs,
while state values depend on earlier states. A future base must not contribute
to an earlier predictive state. Changing the computation order does not permit
changing these information dependencies.}{fig:causality}

For a fixed sequence of coefficients, the update is affine in the incoming
state. The composition of $(a_1,b_1)$ followed by $(a_2,b_2)$ is
\begin{equation}
 (a_2\odot a_1,\ a_2\odot b_1+b_2),\qquad b_t=(1-a_t)\odot v_{i_t}.
\end{equation}
Chapter 7 proved associativity. Here the additional obligation is to establish
that the coefficients are available without first evaluating the state they
will update. Input-only pair lookup satisfies that obligation. It permits an
affine prefix implementation; it does not prove that any particular GPU scan
will be faster than sequential execution for a given size and precision.

Chunk equivalence follows by composing the same transitions in the same order
and carrying the boundary state and previous base. Tests compare every split
of short records and also compare gradients through the split computation.
Detaching a chunk boundary can preserve forward values while changing the
training objective's gradient. Matching output alone is therefore insufficient
for a training-parity claim.

\section{Differentiate the selected operation}
For one scalar channel, hold the selected pair fixed. Write $h'=ah+(1-a)v$,
$g=\sigma(u)$ and $\rho=\sigma(r)$. Then
\begin{align}
 \frac{\partial h'}{\partial h}&=a,\qquad
 \frac{\partial h'}{\partial v}=g(1-\rho),\\
 \frac{\partial h'}{\partial u}&=g(1-g)(1-\rho)(v-h),\\
 \frac{\partial h'}{\partial r}&=g\rho(1-\rho)(h-v).
 \label{eq:derivatives}
\end{align}
To see the signs, increasing write strength moves toward $v$, while increasing
retention moves toward $h$. If $h=v$, both gate derivatives vanish: changing
the mixture cannot change its value. Saturated sigmoids also have small
derivatives. A near-zero gradient may indicate this local geometry rather than
a broken implementation or a useful learned biological switch.

Unselected rows receive no contribution from that event, but a row used several
times accumulates gradient contributions from all those times. This is ordinary
parameter sharing, not one independent parameter per occurrence. The reference
is checked against explicit scalar derivatives, central finite differences and
an independently expanded recurrence. These checks establish arithmetic
consistency; they do not establish that gradient descent will find a good model.

For input-only coefficients, the state-to-state Jacobian across positions is
\begin{equation}
 \frac{\partial h_t}{\partial h_s}
 =\operatorname{diag}\left(\prod_{j=s+1}^{t}a_j\right).
\end{equation}
No coordinate amplification occurs through this path when $a_j\in[0,1]$,
but products can vanish. Candidate and gate parameter gradients include their
own local factors and repeated-use sums. It would be incorrect to extend the
state-path bound to every parameter gradient without analyzing those terms.

\section{A feedback gate changes both the proof and the scheduler}
Suppose a regulator reads the old state, rather than only the input. For one
channel let $g(h)=\sigma(wh)$ and set $f(h)=(1-g(h))h+g(h)v$.
Differentiation gives
\begin{equation}
 f'(h)=1-g(h)+(v-h)w g(h)(1-g(h)).
\end{equation}
At $h=0$, $v=1$ and $w=8$, $f'(0)=1/2+8/4=2.5$.
The output is still a convex combination of $h$ and $v$, yet nearby inputs
can separate locally. Bounded values are not the same property as contraction.
\Plate{04-feedback}{The declared feedback-gated map near zero has slope 2.5,
whereas the identity has slope one. The gate's derivative is part of the
Jacobian. Treating a computed gate as a constant would miss the amplification.}{fig:feedback}
\Listing{feedback_step}

Once coefficients depend on $h_{t-1}$, they cannot generally be prepared as a
token-only affine scan input. Function composition remains associative in the
abstract, but there need not be a closed, fixed-size affine representation of
those composed functions. A claim of parallel execution now needs a different
algorithm or additional structure. This distinction prevents a common category
error: importing a scan theorem after changing the regulator's dependencies.

\section{Gate interpretations may be non-identifiable}
The forward operator uses $g$ and $\rho$ only through $a=1-g(1-\rho)$.
For example, $(g,\rho)=(1/2,1/2)$ and $(3/4,2/3)$ both give $a=3/4$.
They can be represented by finite logits $(0,0)$ and $(\log3,\log2)$.
With identical candidates, every input and incoming state produces the same
forward result. Output fitting alone therefore cannot identify these two
mechanistically named gates separately.

This is structural \Term{non-identifiability}, not merely too little training
data. Separate gate semantics would need additional constraints, independent
observations, or a different architecture. A regularizer can choose one
parameterization but cannot by itself prove a biological interpretation.
Displaying a gate trace is useful for debugging; calling it a discovered
regulatory mechanism requires evidence beyond its name.

\begin{center}\input{\Generated/ablation.tex}\end{center}
The table fixes candidates and changes one logit family at a time. Near-zero
write gates and near-unit retention both make state changes tiny. Neither
intervention is exactly identity at finite logits. These are functional
ablations of an assigned model, not trained accuracy results. A stronger study
would retrain matched alternatives, preserve budgets, repeat seeds and report
both utility and cost with uncertainty.

\section{Place the operator against current research}
Mamba makes selected state-space parameters input dependent and provides a
hardware-aware execution strategy \cite{mamba}. That motivates separating
selection semantics from the implementation used to evaluate them. Our pair
table and convex interpolation are not a Mamba block and inherit none of its
benchmark results. They also make a different representational choice: a
finite pair address is much narrower context than a learned contextual feature.

The 2026 Evo 2 paper describes a multi-hybrid architecture combining several
input-dependent convolution operators and attention \cite{evo2}. Its existence
is a reason to compare complementary mechanisms carefully, not to declare a
non-Transformer direction automatically superior. Evo 2 is an external model,
not this book's Evolutor project. We do not equate a similar name with shared
code, architecture, training data or validation.

An academically useful DOGMA claim should name the task, input and label
semantics, hypothesis, baselines and failure criterion. An industrially useful
claim additionally needs stable model/state versions, record isolation, memory
and latency measurements on specified hardware, failure recovery and reproducible
artifacts. Neither standard is met by synthetic arithmetic alone. The present
operator supplies a small part of the correctness foundation on which those
experiments could be built.

\section{A falsifiable evaluation contract}
For a benign sequence-prediction or regulatory-label study, first declare
whether strand reversal preserves a scalar label or permutes strand-specific
labels. Our directional recurrence is not automatically reverse-complement
equivariant. The full-record averaging construction in Chapter 18 can restore
a specified symmetry, but it changes access and streaming assumptions.

Then separate homologous or duplicated records before constructing train,
validation and held-out test sets; otherwise related sequences can produce
misleading generalization. Choose the split rule for the scientific question:
held-out loci, individuals, chromosomes or species answer different questions.
Fit normalization and selection rules on training data only, and keep model
selection separate from final reporting. This is a proposed evaluation contract,
not a claim that the reference has passed a genomic benchmark.

Compare pair histograms, fixed-retention recurrence, input-conditioned recurrence
and a suitably budgeted sequence baseline. Record parameter count, valid tokens,
training compute, peak memory, preprocessing and uncertainty across independent
runs. Measure gate interventions separately from retrained ablations. Report
failure examples, including long gaps, changed base composition, record boundaries
and reverse-complement transformations. A frontier research program needs
experiments that can reject its favored explanation.

The reference stores $48D$ scalar parameters and $D$ recurrent values, plus
the previous-base field in its caller. Forward arithmetic is $O(nD)$ after
address compilation. Keeping every state for this teaching trace costs $O(nD)$;
fixed recurrent width is not constant training memory. Dense Python validation,
small gathered operations and autograd overhead also matter to runtime.
An operation count is not measured throughput, and no optimized kernel is
delivered by this chapter.

\section{Exercises with worked answers}
\begin{enumerate}
\item With $g=\rho=1/2$, calculate $a$. \textbf{Answer:} $3/4$; the candidate
weight is $1/4$, not the write gate $1/2$.
\item Update $h=0.2$ toward $v=1$. \textbf{Answer:} $0.75(0.2)+0.25=0.4$.
\item Why do AACA and ACAA differ here? \textbf{Answer:} Their candidate
orders differ. The final values are $-1/16$ and $-3/64$, despite equal pair counts.
\item What must cross a raw-sequence chunk boundary? \textbf{Answer:} Recurrent
state and previous valid base, with the same tables and record semantics.
\item A detached boundary matches forward values. Is training parity proved?
\textbf{Answer:} No; detachment removes earlier gradient paths. Check gradients too.
\item Differentiate with respect to $u$ at $h=0,v=1,g=\rho=1/2$.
\textbf{Answer:} $(1/2)(1/2)(1/2)=1/8$.
\item Differentiate with respect to $r$ in the same case. \textbf{Answer:}
$(1/2)(1/2)(1/2)(-1)=-1/8$: retention suppresses the new value.
\item Show non-identifiability using finite logits. \textbf{Answer:}
$(u,r)=(0,0)$ and $(\log3,\log2)$ both produce $a=3/4$.
\item Why does the feedback example defeat a contraction claim?
\textbf{Answer:} Its total derivative at zero is 2.5, including the gate path;
the direct retention coefficient alone is only 0.5.
\item Does every nonlinear recurrence support the affine scan? \textbf{Answer:}
No. A state-dependent regulator generally invalidates token-only coefficient
preparation and closure of the fixed-size affine representation.
\item Count parameters for $D=8$. \textbf{Answer:} $3\cdot16\cdot8=384$,
excluding an output head and any richer context encoder.
\item Design a test of the claim that selective gates improve genomic utility.
\textbf{Rubric:} Fix task and label semantics, leakage-resistant splits,
matched baselines and resource budgets; compare trained controls and interventions,
multiple seeds, uncertainty and out-of-distribution failures. Gate visualization
or one synthetic separation is insufficient evidence.
\end{enumerate}

\section{Vocabulary and next question}
\Term{Input-conditioned coefficient}: an update parameter computed from the
available input, not from a yet-unknown recurrent value.
\Term{State-path Jacobian}: sensitivity propagated through recurrent states.
\Term{Gradient parity}: agreement of derivative semantics, beyond forward values.
\Term{Functional ablation}: changing an operation while holding a declared
parameter set fixed. Chapter 20 asks how to give different memory components
different locality and timescales without hiding their costs or ownership.
